\part{Questions and interface reference}


\section{Frequently asked questions}
\textbf{Why does a long label overflow?} Outside labels need room in addition to
the plot. Increase label width, shorten a unit, or use text to give a shorter label.
Gradtable measures the widest label per column; increase the total width or reduce
padding if its reserved labels leave insufficient graphic space. Unbreakable text
still needs explicit break opportunities.

\textbf{How do I adjust row height?} In gradtable use row stretch and row padding;
for a standalone graphic use row padding. In a standard tabular, arraystretch also
changes row height. References and change annotations may need additional headroom.

\textbf{Why do equal values have different lengths?} Use the same min, max and
graphic width within a comparison column. Spark auto ranges rescale each series;
use fixed for magnitude comparisons. Formatting does not change geometry.

\textbf{Is NA the same as zero?} No. Zero is an observed value and participates in
ranking. NA is missing. A spark keeps its missing position and breaks the line.
A missing stack contribution cannot safely be treated as zero.

\textbf{Why did a table not break across pages?} Normal gradtable is indivisible.
Set \key{long=true} in single-column body text, outside a float or minipage. Individual
rows cannot split. A group heading stays with its first row, but its name is not
repeated on subsequent pages.

\textbf{Why are two identical-looking labels ranked differently?} Ranking uses raw
values, not rounded labels. Increase precision. Ties mark all equal raw values;
second best means the next distinct value in the chosen table or group scope.
